% === Begin Label Data ===
% ID: 605
% 难度星级: 
% 标签: 
% 备注: 
% 组卷引用次数: 0
% === End  Label Data ===

\begin{problem}{2005}{G}{湖北卷（理）}{14}{排列组合}
$ \left(\dfrac{x}{2}+\dfrac{1}{x}+\sqrt{2}\right)^5 $ 的展开式中的常数项为 \underline{\hspace{4em}}.
\end{problem}

\begin{answer}
$\dfrac{63\sqrt{2}}{2}$
\end{answer}

\begin{solutions}
思路 1：可不妨令 $x>0$，利用 $\left(\dfrac{x}{2}+\dfrac{1}{x}+\sqrt{2}\right)=\dfrac{1}{2}\left(\sqrt{x}+\sqrt{\dfrac{2}{x}}\right)^2$，化三项式为二项式，再用二项式定理写出展开式的通项，从而求出展开式的常数项。

当 $x>0$ 时，$\left(\dfrac{x}{2}+\dfrac{1}{x}+\sqrt{2}\right)^5=\left[\dfrac{1}{2}\left(\sqrt{x}+\sqrt{\dfrac{2}{x}}\right)^2\right]^5=2^{-5}\left(\sqrt{x}+\sqrt{\dfrac{2}{x}}\right)^{10}$.

则 $\left(\sqrt{x}+\sqrt{\dfrac{2}{x}}\right)^{10}$ 的二项展开式的通项为 $T_{r+1}=C_{10}^r\cdot x^{\frac{r}{2}-\frac{10-r}{2}}\cdot 2^{\frac{10-r}{2}}$，

常数项所对应的 $r$ 满足 $\dfrac{r}{2}-\dfrac{10-r}{2}=0$，得 $r=5$.

所以所求常数项为 $2^{-5}\cdot C_{10}^5\cdot 2^{\frac{5}{2}}=\dfrac{63\sqrt{2}}{2}$.

思路 2：也可通过先通分化简后，化三项式为二项式，再用二项式定理求出展开式中的常数项。

$\left(\dfrac{x}{2}+\dfrac{1}{x}+\sqrt{2}\right)^5=\left(\dfrac{x^2+2\sqrt{2}x+2}{2x}\right)^5=\dfrac{[(x+\sqrt{2})^2]^5}{(2x)^5}=\dfrac{(x+\sqrt{2})^{10}}{(2x)^5}$.

对于二项式 $(x+\sqrt{2})^{10}$，其通项 $T_{r+1}=C_{10}^r\cdot x^{10-r}\cdot 2^{\frac{r}{2}}$.要得到原展开式的常数项，则需 $10-r=5$，即 $r=5$.

所以所求常数项为 $\dfrac{C_{10}^5\cdot 2^{\frac{5}{2}}}{2^5}=\dfrac{63\sqrt{2}}{2}$.

思路 3：若未能直接看出上述化三项式为二项式的特征，则可将三项式中的前两项先组合起来，再用二项式定理直接展开，解法如下：

$\left(\dfrac{x}{2}+\dfrac{1}{x}+\sqrt{2}\right)^5=\left[\left(\dfrac{x}{2}+\dfrac{1}{x}\right)+\sqrt{2}\right]^5=\sum_{k=0}^{5}C_5^k\left(\dfrac{x}{2}+\dfrac{1}{x}\right)^k\cdot 2^{\frac{5-k}{2}}$，

注意其中常数项只出现在 $k=0,2,4$ 中，合并可得常数项为

$2^{\frac{5}{2}}+C_5^2\cdot 2^{\frac{3}{2}}+C_5^4C_4^2\left(\dfrac{1}{2}\right)^2\cdot 2^{\frac{1}{2}}=\dfrac{63\sqrt{2}}{2}$.
\end{solutions}